\setcounter{subsection}{0}
\section{DESCRIPCIÓN GENERAL}

La introducción mostró que las familias pierden tiempo y dinero al traducir la lista del colegio en
productos y que ninguna herramienta disponible lo resolvía de punta a punta. Este capítulo convierte ese
hallazgo en una pregunta de investigación, presenta la solución (Tukit) y sus objetivos, y explica con qué
metodología se trabajó: una híbrida en la que cada fase respondió a un objetivo y cada método tuvo una
función propia.\par

\subsection{Oportunidad, Problema}

\subsubsection{Contexto del Problema}

El comercio electrónico en Colombia creció de forma sostenida: en 2025 las ventas en línea alcanzaron
COP 105 billones, un 5,9\,\% más que el año anterior \citep{noauthor_informe_nodate-1}. Ese crecimiento no
cambió la manera de adquirir los útiles escolares, una tarea anual que, solo en Bogotá, involucra a
1.171.390 estudiantes \citep{noauthor_estadisticas_nodate}. Según FENALCO, el 35\,\% de las familias
destina entre \$300.000 y \$500.000 por niño \citep{noauthor_resultados_2025}, y el 70\,\% no planea la
compra con anticipación \citep{noauthor_resultados_2024}. A esa presión se suma el esfuerzo mental de
buscar cada artículo por separado y compararlo a mano. La teoría de la carga cognitiva establece que, cuando
la complejidad de una tarea supera la memoria de trabajo, la calidad de las decisiones disminuye
\citep{sweller_cognitive_2011}, y Iyengar y Lepper encontraron que, al pasar de 6 a 24 alternativas, la tasa
de compra bajó del 30\,\% al 3\,\% \citep{iyengar_when_2000}.\par

Los actores que lideran este segmento en Colombia trasladaron su modelo físico al canal digital sin
modificarlo: publican un catálogo y delegan en el cliente la búsqueda, la comparación y la validación
\citep{panamericana_cotizador_2025,primavera_armatulista_2026,mercadolibre_kits_2025,exito_listas_2025}.
TeacherLists, en el ámbito internacional, convierte las listas de los salones en carritos prellenados
\citep{teacherlists_retail_2026}, pero no interpreta listas impresas o fotografiadas ni reconoce el
vocabulario local. Se identificó así la ausencia de un actor que convierta el modelo reactivo de compra en
uno proactivo y asistido, y esa brecha constituyó la oportunidad del proyecto.\par

\subsubsection{Formulación del Problema}

Los canales de venta de útiles son abundantes, pero ninguno traduce de manera automática el texto libre y
heterogéneo de una lista escolar en un pedido correcto y entregable. Esa traducción exige determinar qué
significa cada renglón dentro de un vocabulario que nunca ha sido estandarizado, y establecer qué
combinación de productos equilibra el costo, la marca preferida y la disponibilidad. Por eso el problema se
descompone en tres dimensiones que se resuelven en orden. La percepción convierte una lista recibida en PDF,
fotografía o texto libre en una representación estructurada. La normalización asocia cada renglón, con
frecuencia vago o mal escrito, con el producto del catálogo que le corresponde. La optimización construye
kits que consideran a la vez el costo, las preferencias de la familia y la calidad, y presenta pocas
alternativas diferenciadas.\par

La fricción, entendida como el esfuerzo innecesario en la compra, se analiza con el enfoque de completación
de tareas de Sismeiro y Bucklin \citep{sismeiro_modeling_2004}: cada renglón que el acudiente debe buscar,
comparar y validar por su cuenta es una tarea adicional y un punto donde puede abandonar la compra. Una
solución tecnológica es pertinente porque los agentes de recomendación reducen el esfuerzo de búsqueda y
mejoran la decisión \citep{haubl_consumer_2000}, y porque la madurez reciente de los modelos multimodales
\citep{shen_ocr_2026}, la búsqueda vectorial \citep{karpukhin_dense_2020} y la optimización multiobjetivo
\citep{deb_fast_2002} permite abordar las tres dimensiones en un mismo sistema. De ahí surge la pregunta de
investigación.\par

\begin{quote}
\itshape ¿Cómo construir kits de útiles escolares personalizados y alineados con los
requerimientos específicos de cada institución educativa, considerando el perfil
socioeconómico y las preferencias de la familia, de manera que se reduzca la carga cognitiva
del proceso de compra?
\end{quote}

\subsubsection{Propuesta de Solución}

La propuesta pertenece a la ingeniería de software, con componentes de inteligencia artificial y de
optimización multiobjetivo. Tukit es un sistema de recomendación con un modelo B2B2C
(\emph{business-to-business-to-consumer}): el equipo de Tukit carga y publica la lista oficial de cada
colegio, y el acudiente elige el kit y hace el pedido. El sistema recibe la lista en imagen, PDF o texto
libre y la convierte en kits completos mediante las tres etapas ya descritas, con revisión humana de los
renglones sin certeza suficiente. El perfil socioeconómico de la pregunta se operacionalizó como el
presupuesto y el tope de gasto que declara cada acudiente.\par

Intervienen dos actores con un canal propio. El equipo de Tukit opera una aplicación de administración, en
la que revisa las listas y atiende la cola de renglones dudosos. El acudiente accede por una aplicación web
progresiva, selecciona el colegio y el grado de cada hijo, o carga su lista, e indica presupuesto y
preferencias. Recibe hasta cuatro kits: el Ahorrador (menor costo) y el Premium (mejores marcas) salen
siempre; el Fan, con los personajes pedidos, sale solo si el acudiente pidió uno; y A tu Medida, la mejor
marca posible dentro del monto, sale cuando fija un tope. Después puede reemplazar productos, elegir el kit
y pagar, en modo simulado en esta versión. Un canal para que el colegio publique sus propias listas quedó
como trabajo futuro.\par

\subsubsection{Justificación de la Solución}

El sistema es pertinente porque simplifica un proceso que hoy recae por completo sobre las familias.
Presentar pocas alternativas ya armadas facilita la decisión, porque el exceso de opciones dificulta la
elección \citep{iyengar_when_2000}, y a la vez beneficia a cada actor: el acudiente evita recorrer varios
establecimientos y puede ajustar el kit; el colegio logra que las familias compren a partir de su lista
oficial; y el equipo de Tukit concentra su revisión en los renglones inciertos, de acuerdo con las pautas que
recomiendan indicar lo que el sistema desconoce y facilitar la corrección \citep{amershi_guidelines_2019}.\par

Técnicamente es viable porque las técnicas tienen respaldo empírico: los modelos multimodales leen documentos
con un desempeño comparable al de los métodos tradicionales \citep{shen_ocr_2026} y NSGA-II es una referencia
para objetivos en conflicto \citep{deb_fast_2002}. En la prueba ESC-D4, el armado de kits tuvo un percentil 95
de 0,31~s frente al límite de 8~s (Sección~\ref{sec:optimizacion}), medición que cubre la optimización y no la
lectura de la lista. La arquitectura de microservicios permitiría que cada componente crezca según su demanda
\citep{newman_building_2021}, aunque el despliegue fue local y no se ejercitó el escalado independiente.
El sistema usa herramientas de código abierto, como PostgreSQL, RabbitMQ, Redis y pymoo
\citep{blank_pymoo_2020}, y paga los modelos de inteligencia artificial por uso; como cada lista oficial se
procesa una vez y sirve a todas las familias del grado, el costo se reparte entre muchos usuarios.\par

Para estimar la viabilidad económica se midió el valor de los kits canónicos que el sistema publicó para los
trece grados del Colegio Agustiniano Ciudad Salitre, y el kit Ahorrador costó en promedio \$609.343 en útiles
escolares, sin contar los libros del plan lector. Sobre ese valor se supuso una comisión del 15\,\%,
equivalente a \$91.401 por kit; al descontar la pasarela de pagos Wompi con IVA (\$20.049, según su tarifa
publicada), la inteligencia artificial (\$1.000), la entrega a domicilio en Bogotá (\$33.000) y la adquisición
de clientes (\$8.000), cada kit dejó una contribución de \$29.353. Con la infraestructura del piloto, de
\$1.200.000 al año, el punto de equilibrio fue de 41 kits, cerca del 3\,\% de las 1.340 familias estimadas del
colegio, y con la infraestructura escalable, de \$15.000.000 al año, fue de 512 kits (cálculo reproducible con
el script \texttt{docs/negocio/punto\_equilibrio.py}). La comisión, la logística, la adquisición y la conversión
de las familias fueron supuestos del equipo que la operación real debía confirmar, y con una comisión del
10\,\% la contribución por kit resultó negativa.\par

\subsection{Descripción del Proyecto}

\subsubsection{Objetivo General}

Diseñar una arquitectura de sistema de recomendación multiobjetivo que asista la
construcción de kits escolares, midiendo su efecto sobre la carga cognitiva percibida, la
usabilidad del sistema y la satisfacción del usuario con padres del sector Salitre de
Bogotá durante 2026.

\subsubsection{Objetivos Específicos}
\label{sec:objetivos}

Cada objetivo corresponde a una fase del proyecto, de modo que el resultado de la fase sirva de evidencia de
su cumplimiento.

\begin{itemize}
    \item \textbf{OE1.} Caracterizar el dominio escolar a través del análisis del negocio
    y el levantamiento de requerimientos con los actores clave, identificando las
    restricciones, necesidades y criterios de éxito del sistema, y estableciendo las
    bases conceptuales y de datos necesarias para el diseño y desarrollo del sistema de
    recomendación.
    \item \textbf{OE2.} Diseñar la arquitectura del sistema de recomendación, definiendo
    los componentes principales, las decisiones de diseño y los atributos de calidad
    mediante la evaluación ATAM \citep{bass_software_2021}, y documentando la solución en
    el Documento de Diseño de Software (SDD) que sirva como guía de referencia para las
    fases de implementación posteriores.
    \item \textbf{OE3.} Desarrollar los módulos de percepción y normalización semántica
    del sistema, modelando la extracción de información desde fuentes heterogéneas (PDF,
    imágenes y texto libre) e implementando una búsqueda híbrida, léxica y semántica, para
    la correspondencia entre ítems de la lista institucional y el catálogo de productos.
    \item \textbf{OE4.} Implementar la generación de recomendaciones personalizadas
    mediante un enfoque multiobjetivo que logre un equilibrio técnico entre el costo, la
    fidelidad de marca y la disponibilidad, permitiendo la construcción de configuraciones
    diferenciadas que satisfagan las diversas restricciones de los perfiles de usuario
    identificados y traduciendo las preferencias en lenguaje natural a restricciones
    estructuradas (NL2C), integradas en el sistema con una interfaz web progresiva.
    \item \textbf{OE5.} Evaluar el sistema mediante pruebas funcionales, no funcionales,
    de carga y de usabilidad con usuarios reales, verificando el correcto funcionamiento
    de cada componente del sistema y midiendo la experiencia del usuario en términos de
    satisfacción y reducción de la carga cognitiva percibida
    \citep{sweller_cognitive_2011} frente al proceso tradicional
    de compra.
\end{itemize}

\subsection{Marco metodológico, entregables, estándares y justificación}

El proyecto siguió una metodología híbrida contenedora, llamada así porque reúne varios marcos y le asigna a
cada uno una función distinta. SCRUM \citep{schwaber_scrum_2020} organizó la construcción en ciclos de dos
semanas, con un tablero Kanban para el estado de cada tarea. Design Thinking \citep{brown_design_2008} guió
la empatía y la definición del OE1, es decir, las entrevistas, la encuesta, la propuesta de valor y el
modelo de negocio. RUP \citep{kruchten_rup_2003} aportó la gestión de riesgos y la trazabilidad de cada
requerimiento hasta su componente y su prueba. ATAM \citep{bass_software_2021} permitió elegir la
arquitectura del OE2 entre tres alternativas. CRISP-DM \citep{chapman_crisp-dm_2000,wirth_crisp-dm_2000}
fue el proceso de los componentes que dependen de datos (OE3 y OE4), y para el OE5 se usó ISO/IEC~25010
\citep{iso_25010_2011} con un protocolo de validación registrado de antemano. Las normas IEEE 830, 1016 y 1058
definieron la estructura de los documentos de ingeniería: requerimientos, diseño y plan del proyecto.\par

El trabajo se dividió en cinco fases, una por objetivo específico. La Tabla~\ref{tab:fases} muestra, para
cada fase, el objetivo, la etapa de CRISP-DM que le corresponde, los métodos y lo que produjo; conviene
leerla de arriba abajo, porque sigue el orden en que se trabajó y el de los capítulos siguientes.\par

\begin{table}[H]
\centering\footnotesize
\setlength{\tabcolsep}{4pt}
\renewcommand{\arraystretch}{1.1}
\caption{Fases del proyecto, su relación con CRISP-DM y sus métodos.}
\label{tab:fases}
\begin{tabularx}{\textwidth}{@{}p{2.3cm}p{0.8cm}p{2.4cm}Xp{2.9cm}@{}}
\toprule
Fase & Obj. & Etapa de CRISP-DM & Métodos y técnicas & Resultado \\
\midrule
Caracterización del problema & OE1 & Comprensión del negocio &
Entrevistas, mapas de empatía, encuesta a acudientes, lienzo del modelo de negocio
\citep{osterwalder_business_2010}, FODA y recolección de listas &
Requerimientos, modelo de negocio y catálogo de conceptos \\
Diseño de la arquitectura & OE2 & No aplica (ingeniería de software) &
Atributos de calidad \citep{richards_fundamentals_2025,iso_25010_2011}, árbol de
utilidad, ATAM sobre tres alternativas \citep{bass_software_2021} y vistas 4+1
\citep{kruchten_41_1995} &
Arquitectura y documento de diseño \\
Percepción y normalización & OE3 & Datos, preparación, modelado y evaluación &
Corpus de oro anotado por una persona; kappa de Cohen \citep{cohen_coefficient_1960} entre
las reglas y un modelo de lenguaje; comparación de motores de lectura y de estrategias de
búsqueda; análisis del umbral con revisión humana &
Módulos de percepción y normalización \\
Optimización e integración & OE4 & Modelado, evaluación y despliegue &
Formulación multiobjetivo, NSGA-II \citep{deb_fast_2002}, búsqueda exacta, utilidad
multiatributo (MAUT), traducción de preferencias (NL2C) e integración en sprints &
Optimizador, aplicación web y aplicación de administración \\
Validación & OE5 & Evaluación &
Pruebas funcionales, no funcionales y de carga; SUS \citep{brooke_sus_1996}, NPS
\citep{reichheld_one_2003} y verbalización concurrente \citep{lazar_research_2017} &
Informe de pruebas y evaluación con usuarios \\
\bottomrule
\end{tabularx}
\end{table}

\subsubsection{Extensión de CRISP-DM a los componentes que dependen de datos}
\label{sec:crispdm-extension}

CRISP-DM define seis etapas: comprensión del negocio, comprensión de los datos, preparación, modelado,
evaluación y despliegue. En la Fase~1 se aplicó solo la primera. Desde la Fase~3 se extendió a todos los
componentes cuyo comportamiento depende de datos: la percepción y la normalización del OE3, y la
optimización, la utilidad multiatributo (MAUT) y la traducción de preferencias a restricciones (NL2C) del
OE4. La Tabla~\ref{tab:crispdm-oe} indica, por objetivo, qué se hizo en cada etapa, qué evidencia produjo y
en qué sección de esta memoria se presenta.\par

\begin{table}[H]
\centering\footnotesize
\setlength{\tabcolsep}{3pt}
\renewcommand{\arraystretch}{1.1}
\caption{Extensión de CRISP-DM por objetivo: etapa, actividad, evidencia y sección.}
\label{tab:crispdm-oe}
\begin{tabularx}{\textwidth}{@{}p{0.9cm}p{2.4cm}XXp{1.9cm}@{}}
\toprule
OE & Etapa & Actividad & Evidencia & Sección \\
\midrule
OE1 & Comprensión del negocio & Entrevistas, encuesta a acudientes (N=65 recibidas, n=60
analizables) y levantamiento de listas de 17 colegios & Dossier de negocio y SRS &
Cap.~III y IV \\
\midrule
OE3 & Comprensión de los datos & Caracterización del corpus de listas (4.719 ítems) y
armado del conjunto de oro CGS-10 (diez documentos, 1.293 ítems) & Corpus y conjunto de oro &
Sec.~\ref{sec:percepcion} \\
OE3 & Preparación & Limpieza del texto, corrección ortográfica con el léxico del catálogo, núcleo
en singular y separación de atributos & Código de normalización &
Sec.~\ref{sec:normalizacion} \\
OE3 & Modelado & Cinco enfoques de lectura; cuatro estrategias de recuperación; resolución
léxica con aporte del coseno & Bancos de los prototipos &
Sec.~\ref{sec:percepcion-dev} y \ref{sec:normalizacion} \\
OE3 & Evaluación & F1 de la lectura; análisis del umbral de certeza $\tau$, declarado sin calibrar; revisión
humana de los renglones dudosos & Tablas de resultados &
Sec.~\ref{sec:percepcion} y \ref{sec:recuperacion} \\
OE3 & Despliegue & Servicios de inteligencia y de catálogo; cola de revisión del administrador &
Sistema construido & Sec.~\ref{sec:pipeline} \\
\midrule
OE4 & Comprensión de los datos & Encuesta de preferencias (pregunta P21, 60 respuestas) y 174
listas del corpus & Dossier de negocio y banco de listas & Sec.~\ref{sec:optimizacion-dev} \\
OE4 & Preparación & Líneas comprables, candidatos por renglón y curaduría del catálogo &
Auditorías de los kits & Sec.~\ref{sec:optimizacion-dev} \\
OE4 & Modelado & NSGA-II con cuatro objetivos, frente exacto, MAUT y NL2C &
Código y banco de calibración & Sec.~\ref{sec:optimizacion-cal} y \ref{sec:nl2c} \\
OE4 & Evaluación & Comparación con el óptimo exacto, hipervolumen y cuatro auditorías de kits &
Tablas de resultados & Sec.~\ref{sec:optimizacion} \\
OE4 & Despliegue & Aplicación web, aplicación de administración y pago & Sistema
construido & Sec.~\ref{sec:web} \\
\bottomrule
\end{tabularx}
\end{table}

\subsubsection{Ciclos de calibración con revisión humana}
\label{sec:ciclos-calibracion}

Las etapas de CRISP-DM se recorrieron en varios ciclos, y en percepción, normalización y optimización cada
ciclo incluyó una revisión humana. Primero se ejecutó el componente sobre una muestra; segundo, una persona
revisó los resultados (el administrador en la cola de renglones dudosos, o el equipo al auditar los kits
renglón por renglón); tercero, esa revisión produjo etiquetas o correcciones; y cuarto, el equipo ajustó
los datos, las reglas o las constantes y repitió la evaluación con la misma métrica. Estos ciclos no son
aprendizaje por refuerzo, porque no hay señal de recompensa ni una política que se actualice sola: cada
cambio queda registrado con su evidencia en la Tabla~\ref{tab:constantes} (Anexo~\ref{anx:constantes}).
El diseño sigue las pautas de interacción humano-IA de Amershi et al. \citep{amershi_guidelines_2019}. El
Capítulo~VI describe cada ciclo junto a su componente, y el despliegue con colegios y familias se presenta
como trabajo futuro en el Capítulo~VIII. El Capítulo~III abre el contexto que sustenta estas fases.\par
